\nonstopmode{}
\documentclass[a4paper]{book}
\usepackage[times,inconsolata,hyper]{Rd}
\usepackage{makeidx}
\makeatletter\@ifl@t@r\fmtversion{2018/04/01}{}{\usepackage[utf8]{inputenc}}\makeatother
% \usepackage{graphicx} % @USE GRAPHICX@
\makeindex{}
\begin{document}
\chapter*{}
\begin{center}
{\textbf{\huge Package `spectrakit'}}
\par\bigskip{\large \today}
\end{center}
\ifthenelse{\boolean{Rd@use@hyper}}{\hypersetup{pdftitle = {spectrakit: Spectral Data Handling and Visualization}}}{}
\ifthenelse{\boolean{Rd@use@hyper}}{\hypersetup{pdfauthor = {Gianluca Pastorelli}}}{}
\begin{description}
\raggedright{}
\item[Type]\AsIs{Package}
\item[Title]\AsIs{Spectral Data Handling and Visualization}
\item[Version]\AsIs{0.1.2}
\item[Description]\AsIs{Provides functions to combine, normalize and visualize spectral data, perform principal component analysis (PCA), and assemble customizable image grids suitable for publication-quality scientific figures.}
\item[License]\AsIs{MIT + file LICENSE}
\item[Encoding]\AsIs{UTF-8}
\item[Imports]\AsIs{dplyr, readr, ggplot2, ggrepel, tibble, purrr, rlang, magick,
glue, data.table, stats, grid}
\item[Suggests]\AsIs{RColorBrewer}
\item[RoxygenNote]\AsIs{7.3.3}
\item[NeedsCompilation]\AsIs{no}
\item[Author]\AsIs{Gianluca Pastorelli [aut, cre] (ORCID:
<}\url{https://orcid.org/0000-0001-6926-1952}\AsIs{>)}
\item[Maintainer]\AsIs{Gianluca Pastorelli }\email{gianluca.pastorelli@gmail.com}\AsIs{}
\end{description}
\Rdcontents{Contents}
\HeaderA{combineSpectra}{Combine and Normalize Spectral Data from Multiple Files}{combineSpectra}
%
\begin{Description}
Reads spectral data from multiple files in a folder, merges them by a common column (e.g., wavelength),
optionally filters by a range, applies normalization and returns the data in either row-wise or column-wise format.
\end{Description}
%
\begin{Usage}
\begin{verbatim}
combineSpectra(
  folder = ".",
  file_type = "csv",
  sep = ",",
  header = TRUE,
  common_col_pos = 1,
  data_col_pos = 2,
  range = NULL,
  normalization = c("none", "simple", "min-max", "z-score"),
  orientation = c("columns", "rows")
)
\end{verbatim}
\end{Usage}
%
\begin{Arguments}
\begin{ldescription}
\item[\code{folder}] Character. Path to the folder containing spectral files. Default is working directory (`"."`).

\item[\code{file\_type}] Character. File extension (without dot) to search for. Default is `"csv"`.

\item[\code{sep}] Character. Delimiter for file columns. Use `","` for comma-separated (default) or `"\bsl{}t"` for tab-delimited files.

\item[\code{header}] Logical. Whether the files contain a header row. Default is `TRUE`.

\item[\code{common\_col\_pos}] Integer. Column position for the common variable (e.g., wavelength). Defaults to `1`.

\item[\code{data\_col\_pos}] Integer. Column position for the spectral intensity values. Defaults to `2`.

\item[\code{range}] Numeric vector of length 2. Optional range filter for the common column (e.g., wavelength limits). Defaults to `NULL` (no filtering).

\item[\code{normalization}] Character. Normalization method to apply to spectra. One of `"none"`, `"simple"` (divide by max), `"min-max"`, or `"z-score"`. Default is `"none"`.

\item[\code{orientation}] Character. Output orientation. Use `"columns"` (default) to keep each spectrum as a column, or `"rows"` to transpose so each spectrum is a row.
\end{ldescription}
\end{Arguments}
%
\begin{Value}
A `tibble` that can be exported as, for example, a CSV file. Each spectrum is either a column (default) or row, depending on `orientation`; a sum or mean spectrum can optionally be computed using `rowSums()` or `colSums()`, or `rowMeans()` or `colMeans()`, respectively. The common column (e.g., wavelength) is retained.
\end{Value}
%
\begin{Examples}
\begin{ExampleCode}
# Create a temporary directory for mock CSV files
tmp_dir <- tempdir()

# Define file paths
tmp1 <- file.path(tmp_dir, "file1.csv")
tmp2 <- file.path(tmp_dir, "file2.csv")

# Write two mock CSV files in the temporary folder
write.csv(data.frame(ID = c("A", "B", "C"), val = c(1, 2, 3)), tmp1, row.names = FALSE)
write.csv(data.frame(ID = c("A", "B", "C"), val = c(4, 5, 6)), tmp2, row.names = FALSE)

# Merge the CSV files in the temporary folder, normalize with z-score, and return transposed
result <- combineSpectra(
  folder = tmp_dir,
  file_type = "csv",
  sep = ",",
  common_col_pos = 1,
  data_col_pos = 2,
  normalization = "z-score",
  orientation = "rows"
)

\end{ExampleCode}
\end{Examples}
\HeaderA{makeComposite}{Create a labeled image grid}{makeComposite}
%
\begin{Description}
Generates a composite image grid with customizable layout, labels and resizing options. Suitable for spectra and other image types.
\end{Description}
%
\begin{Usage}
\begin{verbatim}
makeComposite(
  folder = ".",
  custom_order = NULL,
  rows = NULL,
  cols = NULL,
  spacing = 15,
  resize_mode = c("none", "fit", "fill", "width", "height", "both"),
  labels = list(),
  label_settings = list(),
  background_color = "white",
  desired_width = 15,
  width_unit = "cm",
  ppi = 300,
  output_format = "tiff",
  output_folder = NULL
)
\end{verbatim}
\end{Usage}
%
\begin{Arguments}
\begin{ldescription}
\item[\code{folder}] Character. Path to the folder containing images. Default is working directory (`"."`).

\item[\code{custom\_order}] Character vector. Set of filenames (use NA for blank slots).

\item[\code{rows}] Integer. Number of rows in the grid.

\item[\code{cols}] Integer. Number of columns in the grid.

\item[\code{spacing}] Integer. Spacing (in pixels) between tiles. Default is `15`

\item[\code{resize\_mode}] Character. Method to resize panels in the composite. Options are:
\begin{description}

\item[`"none"`] Keep each panel at its original size.
\item[`"fit"`] Scale each panel to fit within the smallest image dimensions, preserving aspect ratio.
\item[`"fill"`] Scale and crop each panel to exactly fill the smallest dimensions.
\item[`"width"`] Resize each panel to the minimum width, keeping the original height.
\item[`"height"`] Resize each panel to the minimum height, keeping the original width.
\item[`"both"`] Force panels to the exact width and height, which may distort aspect ratio.

\end{description}


\item[\code{labels}] List of up to 4 character vectors. Each vector corresponds to one label layer and must be the same length as the number of non-NA images. Use empty strings "" or NULL entries to omit specific labels.

\item[\code{label\_settings}] List of named lists. Each named list specifies settings for a label layer. Options include:
\begin{description}

\item[`size`] Font size (e.g., 100).
\item[`color`] Font color.
\item[`font`] Font family (e.g., `"Arial"`).
\item[`boxcolor`] Background color behind text, or `NA` for none.
\item[`location`] Offset from the gravity anchor (e.g., `"+10+10"`).
\item[`gravity`] Placement anchor for the label (e.g., `"northwest"`).
\item[`weight`] Font weight (e.g., `400` = normal, `700` = bold).

\end{description}


\item[\code{background\_color}] Character. Background color used for blank tiles and borders. Use `"none"` for transparency. Default is `"white"`.

\item[\code{desired\_width}] Numeric. Desired width of final image (in cm or px). Default is `15`

\item[\code{width\_unit}] Character. Either "cm" or "px". Default is `"cm"`

\item[\code{ppi}] Numeric. Resolution (pixels per inch) for output file. Default is `300`

\item[\code{output\_format}] Character. File format for saving plots. Examples: `"tiff"`, `"png"`, `"pdf"`. Default is `"tiff"`.

\item[\code{output\_folder}] Character. Path to folder where image is saved. If NULL (default), image is not saved; if `"."`, image is saved in the working directory.
\end{ldescription}
\end{Arguments}
%
\begin{Value}
Saves image composite to a specified output folder. Returns `NULL` (used for side-effects).
\end{Value}
%
\begin{Examples}
\begin{ExampleCode}
library(magick)

tmp_dir <- file.path(tempdir(), "spectrakit_imgs")
dir.create(tmp_dir, showWarnings = FALSE)

# Create and save img1
img1 <- image_blank(100, 100, "white")
img1 <- image_draw(img1)
symbols(50, 50, circles = 30, inches = FALSE, add = TRUE, bg = "red")
dev.off()
img1_path <- file.path(tmp_dir, "img1.png")
image_write(img1, img1_path)

# Create and save img2
img2 <- image_blank(100, 100, "white")
img2 <- image_draw(img2)
rect(20, 20, 80, 80, col = "blue", border = NA)
dev.off()
img2_path <- file.path(tmp_dir, "img2.png")
image_write(img2, img2_path)

# Create composite
makeComposite(
        folder = tmp_dir,
        custom_order = c("img1.png", "img2.png"),
        rows = 1,
        cols = 2,
        labels = list(c("Red Circle", "Blue Rectangle")),
        label_settings = list(
                list(size = 5, font = "Arial", color = "black", boxcolor = "white",
                     gravity = "northwest", location = "+10+10", weight = 400)
        ),
        resize_mode = "none",
        desired_width = 10,
        width_unit = "cm",
        ppi = 300,
        output_format = "png",
        output_folder = tmp_dir
)

\end{ExampleCode}
\end{Examples}
\HeaderA{plotPCA}{Perform PCA and Create a Plot of Scores, Loadings, or Biplot}{plotPCA}
%
\begin{Description}
Computes principal component analysis (PCA) on numeric variables in a dataset
and generates a PC1 vs PC2 visualization (scores, loadings, or biplot).
\end{Description}
%
\begin{Usage}
\begin{verbatim}
plotPCA(
  data,
  color_var = NULL,
  shape_var = NULL,
  plot_type = c("score", "loadings", "biplot"),
  palette = "Dark2",
  show_labels = TRUE,
  ellipses = FALSE,
  ellipse_var = NULL,
  display_names = TRUE,
  legend_title = NULL,
  return_pca = FALSE,
  output_format = "tiff",
  output_folder = NULL
)
\end{verbatim}
\end{Usage}
%
\begin{Arguments}
\begin{ldescription}
\item[\code{data}] A data frame containing numeric variables. Non-numeric columns are ignored.

\item[\code{color\_var}] Optional character. Column name for coloring points by group. Converted to factor internally.

\item[\code{shape\_var}] Optional character. Column name for shaping points by group. Converted to factor internally.

\item[\code{plot\_type}] Character. Type of PCA plot to generate:
\begin{description}

\item[`"score"`] Plot PCA scores (observations).
\item[`"loadings"`] Plot PCA loadings (variables).
\item[`"biplot"`] Combine scores and loadings in a biplot.

\end{description}


\item[\code{palette}] Character or vector. Color palette for groups:
\begin{description}

\item[`"Dark2"`] Use Dark2 palette from RColorBrewer (requires the package).
\item[single color] A single color repeated for all groups.
\item[vector of colors] Custom vector of colors, recycled to match number of groups.

\end{description}


\item[\code{show\_labels}] Logical. Display labels for points (scores) or variables (loadings). Default is TRUE.

\item[\code{ellipses}] Logical. Draw confidence ellipses around groups in score/biplot. Grouping logic for ellipses follows this priority:
\begin{itemize}

\item{} If \code{ellipse\_var} is provided, ellipses are drawn by that variable.
\item{} Else, if \code{color\_var} is provided, ellipses are drawn by color groups.
\item{} Else, if \code{shape\_var} is provided, ellipses are drawn by shape groups.
\item{} If none are provided, no ellipses are drawn.

\end{itemize}

Default is FALSE.

\item[\code{ellipse\_var}] Optional character. Name of the variable used to group ellipses. Takes precedence over all other grouping variables. Converted to factor internally. Default is NULL.

\item[\code{display\_names}] Logical. Show legend if TRUE. Default is TRUE.

\item[\code{legend\_title}] Optional character. Legend title corresponding to `color\_var` or `shape\_var`. Default is NULL.

\item[\code{return\_pca}] Logical. If TRUE, return a list with plot and PCA object. Default is FALSE.

\item[\code{output\_format}] Character. File format for saving plots. Examples: `"tiff"`, `"png"`, `"pdf"`. Default is `"tiff"`.

\item[\code{output\_folder}] Character. Path to folder where plots are saved. If NULL (default), returns a ggplot object (or list with plot and PCA if `return\_pca = TRUE`). If specified, plot is saved automatically (function returns PCA object only if `return\_pca = TRUE`); if `"."`, plot is saved in the working directory.
\end{ldescription}
\end{Arguments}
%
\begin{Value}
A ggplot2 object representing the PCA plot, or a list with `plot` and `pca` if `return\_pca = TRUE`.
\end{Value}
%
\begin{Examples}
\begin{ExampleCode}
plotPCA(
  data = iris,
  color_var = "Species",
  shape_var = "Species",
  plot_type = "biplot",
  palette = "Dark2",
  show_labels = TRUE,
  ellipses = TRUE,
  display_names = TRUE,
  legend_title = "Iris Species"
)

\end{ExampleCode}
\end{Examples}
\HeaderA{plotSpectra}{Plot Spectral Data from Multiple Files}{plotSpectra}
%
\begin{Description}
Reads, normalizes and plots spectral data from files in a folder. Supports multiple plot modes, color palettes, axis customization, annotations and automatic saving of plots to files.
\end{Description}
%
\begin{Usage}
\begin{verbatim}
plotSpectra(
  folder = ".",
  file_type = "csv",
  sep = ",",
  header = TRUE,
  normalization = c("none", "simple", "min-max", "z-score", "area", "vector"),
  x_config = NULL,
  x_reverse = FALSE,
  y_trans = c("linear", "log10", "sqrt"),
  x_label = NULL,
  y_label = NULL,
  line_size = 0.5,
  palette = "black",
  plot_mode = c("individual", "overlapped", "stacked"),
  display_names = FALSE,
  vertical_lines = NULL,
  shaded_ROIs = NULL,
  annotations = NULL,
  output_format = "tiff",
  output_folder = NULL
)
\end{verbatim}
\end{Usage}
%
\begin{Arguments}
\begin{ldescription}
\item[\code{folder}] Character. Path to the folder containing spectral files. Default is working directory (`"."`).

\item[\code{file\_type}] Character. File extension (without dot) to search for. Default is `"csv"`.

\item[\code{sep}] Character. Delimiter for file columns. Use `","` for comma-separated (default) or `"\bsl{}t"` for tab-delimited files.

\item[\code{header}] Logical. Whether the files contain a header row. Default is `TRUE`.

\item[\code{normalization}] Character. Normalization method to apply to y-axis data. Options are:
\begin{description}

\item[`"none"`] No normalization is applied (default).
\item[`"simple"`] Divide by the maximum intensity.
\item[`"min-max"`] Scale intensities to the [0,1] range.
\item[`"z-score"`] Subtract the mean and divide by the standard deviation of intensities.
\item[`"area"`] Divide by the total sum of intensities so the spectrum area = 1.
\item[`"vector"`] Normalize the spectrum as a unit vector by dividing by the square root of the sum of squared intensities (L2 normalization).

\end{description}


\item[\code{x\_config}] Numeric vector of length 3. Specifies x-axis range and breaks: `c(min, max, step)`.

\item[\code{x\_reverse}] Logical. If `TRUE`, reverses the x-axis. Default is `FALSE`.

\item[\code{y\_trans}] Character. Transformation for the y-axis. One of `"linear"`, `"log10"`, or `"sqrt"`. Default is `"linear"`.

\item[\code{x\_label}] Character or expression. Label for the x-axis. Supports mathematical notation via `expression()`.

\item[\code{y\_label}] Character or expression. Label for the y-axis. Supports mathematical notation via `expression()`.

\item[\code{line\_size}] Numeric. Width of the spectral lines. Default is `0.5`.

\item[\code{palette}] Character or vector. Color setting: a single color (e.g., `"black"`), a ColorBrewer palette name (e.g., `"Dark2"`), or a custom color vector.

\item[\code{plot\_mode}] Character. Plotting style. One of `"individual"` (one plot per spectrum), `"overlapped"` (all in one), or `"stacked"` (faceted). Default is `"individual"`.

\item[\code{display\_names}] Logical. If `TRUE`, adds file names as titles to individual spectra or a legend to combined spectra. Default is `FALSE`.

\item[\code{vertical\_lines}] Numeric vector. Adds vertical dashed lines at given x positions.

\item[\code{shaded\_ROIs}] List of numeric vectors. Each vector must have two elements (`xmin`, `xmax`) to define shaded x regions.

\item[\code{annotations}] Data frame with columns `file` (file name without extension), `x`, `y`, and `label`. Adds annotation labels to specific points in spectra.

\item[\code{output\_format}] Character. File format for saving plots. Examples: `"tiff"`, `"png"`, `"pdf"`. Default is `"tiff"`.

\item[\code{output\_folder}] Character. Path to folder where plots are saved. If NULL (default), plots are not saved; if `"."`, plots are saved in the working directory.
\end{ldescription}
\end{Arguments}
%
\begin{Details}
Color settings can support color-blind-friendly palettes from `RColorBrewer`. Use `display.brewer.all(colorblindFriendly = TRUE)` to preview.
\end{Details}
%
\begin{Value}
Saves plots to a specified output folder. Returns `NULL` (used for side-effects).
\end{Value}
%
\begin{Examples}
\begin{ExampleCode}
# Create a temporary directory and write mock spectra files
tmp_dir <- tempdir()
write.csv(data.frame(Energy = 0:30, Counts = rpois(31, lambda = 100)),
          file.path(tmp_dir, "spec1.csv"), row.names = FALSE)
write.csv(data.frame(Energy = 0:30, Counts = rpois(31, lambda = 120)),
          file.path(tmp_dir, "spec2.csv"), row.names = FALSE)

# Plot the mock spectra using various configuration options
plotSpectra(
  folder = tmp_dir,
  file_type = "csv",
  sep = ",",
  normalization = "min-max",
  x_config = c(0, 30, 5),
  x_reverse = FALSE,
  y_trans = "linear",
  x_label = expression(Energy~(keV)),
  y_label = expression(Counts/1000~s),
  line_size = 0.7,
  palette = c("black","red"),
  plot_mode = "overlapped",
  display_names = TRUE,
  vertical_lines = c(10, 20),
  shaded_ROIs = list(c(12, 14), c(18, 22)),
  output_format = "png",
  output_folder = tmp_dir
)

\end{ExampleCode}
\end{Examples}
\printindex{}
\end{document}
